\subsection{Full-Task Simulation Performance}
\begin{outline}
\item \reported{latest run exceeds 50\% success and belongs to the latest method direction.} Its checkpoint is absent from retained reports. Reserve the headline row; obtain exact successes/trials, checkpoint hash, metric definition, split, and selection rule before claiming strict held-out performance.
\item \recorded{older matched development comparison: source 756/2048 (36.91\%), latent-likelihood bounded-residual continuation 843/2048 (41.16\%), added correction-network arm 769/2048 (37.55\%).} These are four repeats on 512 scenes, not 2048 independent scenes.
\item \recorded{control minus source: +4.25 percentage points, paired 95\% CI [+2.10,+6.40], Holm-adjusted $p=0.00030$.} This supports the continued control's gain, not the isolated causal effect of changing likelihood.
\item \recorded{older control strict rates: plate 45.43\%, bowl 36.86\%; physical grasp 76.37\%, held lift 68.80\%, native placement 47.90\%.} Use stage retention to explain what changed.
\item Table I plan: matched baselines and latest method; columns for strict successes/trials, strict rate with uncertainty, plate/bowl rates, native placement, and interaction budget. Separate unmatched historical results from this table.
\item Figure 2 plan: strict validation versus cumulative interaction, with protocol/likelihood changes marked. Distinguish in-run validation from standalone evaluation; no smoothing that hides regression or seed variation.
\end{outline}
